
En
effet, raisonnons par récurrence sur \(r\). Lorsque \(r=0\), c'est clair car \(m\ge0\) et \(F\) est \(n\)-régulier pour tout \(n\). Lorsque \(r\ge1\), on a une suite exacte $0\longrightarrow F(-1)\longrightarrow F\longrightarrow G\longrightarrow0$ où \(G\) est un \((b_1,\ldots,b_r)\)-faisceau. Par 1.7 et l'hypothèse de récurrence on voit que si \(n=P_{r-1}(c_1,\ldots,c_r)\), alors \(n\ge0\) et \(G\) est \(n\)-régulier. Donc, par 1.4 \(F\) est \([n+h^1(F(n-1))]\)-régulier et \(h^q(F(n-1))=0\) pour \(q\ge2\). Donc, \(h^1(F(n-1))=h^0(F(n-1))-\chi(F(n-1))\), ce qui est \(\le \sum_{i=0}^{r}(b_i-a_i)\binom{n-1+i}{i}\) par 1.8 (i) car \(b_i\ge0\). Enfin, \(F\) est \(\left[n+\sum_{i=0}^{r}c_i\binom{n-1+i}{i}\right]\)-régulier en vertu de 1.3 (i); d'où la première assertion. La seconde en résulte grâce à 1.8 (ii) et 1.10.

\textbf{1.12.} Soient \(S\) un schéma noethérien et \(X\) un \(S\)-schéma de type fini. Soit \(\mathcal F\) une \uline{famille de classes de faisceaux cohérents sur les fibres de} \(X/S\); c'est-à-dire, pour tout point \(s\in S\) et toute extension \(K\) de \(k(s)\), on se donne des faisceaux cohérents \(F_K\) sur le schéma algébrique \(X_K\) et on dit qu'un \(F_K\) et un \(F'_{K'}\) sont dans la même classe s'il existe des \(k(s)\)-homomorphismes de \(K,K'\) dans une même extension \(K''\) de \(k(s)\) tels que \(F_{K''}(=F_K\otimes_K K'')\) et \(F'_{K''}\) soient isomorphes sur \(X_{K''}\). Alors la famille \(\mathcal F\) est dite \uline{limitée par le faisceau cohérent} \(F\) \uline{sur} \(X_T=X\times_S T\) où \(T\) est de type fini sur \(S\), si \(\mathcal F\) est contenue dans la famille des classes des faisceaux \(F_{k(t)}\) avec \(t\in T\); on dit que \(\mathcal F\) est \uline{limitée} s'il existe \(T,F\) comme dessus.

Supposons en plus que \(X/S\) soit projectif et muni d'un faisceau inversible (relativement) ample \({O}_X(1)\). Alors \(\mathcal F\) est dite une \((b)\)-\uline{famille}, pour une suite d'entiers \((b)=(b_0,\ldots,b_r)\), si toute classe dans \(\mathcal F\) est représentée par un \(F_K\) avec \(K\) algébriquement clos, qui est un \((b)\)-faisceau.

